\section{Détails}

Dans cette partie, nous allons discuter sur un des paramètres qui jouent sur les détails apportés sur la carte générée.

\subsection{Fréquence}

Un des paramètres du bruit de Perlin est la fréquence dans une grille d'échantillonnage discrète.
Dans notre implémentation, nous avons appelé improprement ce paramètre \texttt{scale}.  
Ce paramètre joue sur le nombre de points dans une cellule conceptuelle.
Plus il y aura de points dans une cellule, plus les valeurs dans la cellulle donnée vont avoir des valeurs proches l'une de l'autre. 

\subsubsection{Cas limites}

Nous étudierons les différents intervalles de la fréquence, \textit{freq} pour l'évaluation de la fonction \texttt{noise (x * freq, y * freq)} où $x,y \in \mathbb{N}$. 

\subsubsection*{Fréquence prend des valeurs entières}

\begin{itemize}
	\item Les points évalués pour la fonction de bruit sont des valeurs entières, donc $\in \mathbb{Z}^2$. Or dans l'algorithme, ces coordonnées correspondent aux intersections des cellules de la grille interne.
	\item Dés lors, les vecteurs de distance calculés sont nuls et le produit scalaire est nul.
	\item L'interpolation entre valeurs nulles donne un résultat nul. Sur toute la grille, notre fonction de bruit est nulle et donne une surface monochrome puisque la fonction de bruit est constante en tout point.
	\item Cela produit une surface de couleur monochrome.
	\item Il s'agit d'un phénomène d'alignement entre la grille d'échantillonnage et la grille interne du bruit.

\end{itemize}

\begin{itemize}
		
\item Dans le cas où \texttt{freq = 0}, nous appliquons la fonction de bruit pour le même point, soit $(0,0)$.
\item Pour  le cas où \texttt{freq = 1}, nous calculons la fonction de bruit pour chaque intersection de cellules distinctes.
\item Alors que dans le cas où \texttt{freq > 1}
	\begin{itemize}
		\item Nous ne calculons pas la fonction de bruit  sur toutes les intersections mais pour toutes les $freq$ intersections de cellules.
	\end{itemize}

\end{itemize}

\subsubsection*{Fréquence prend des valeurs avec une partie décimale non nulle}

\begin{itemize}
	\item Fréquence tend vers 0, en restant strictement différent de 0
		\begin{itemize}
			\item Le domaine de la fonction de bruit est plus grand. Cela se traduit par un plus grand nombre de points pour une seule cellule donnée de la grille interne de Perlin.
			\item Pour des points voisins dans une cellule donnée, leurscoordonnées sont très proches.
			\item Les valeurs de bruit pour les différents points dans une cellules données varient très peu puisque leurs vecteurs de distance sont proches en valeurs pour les mêmes vecteurs de gradient.
			\item Dés lors, la variation de la fonction de bruit devient extrêmement faible. Ce qui donne visuellement une surface presque monochrome, proche du cas où \texttt{freq = 0}.
		\end{itemize}

	\item Fréquence prend des valeurs non entières : $freq \in \mathbb{Q}^+ \backslash \mathbb{N}$ 
	\begin{itemize}
		\item Pour les points considérés en paramètres de \texttt{noise (x * freq, y * freq)}, la fonction de bruit s'applique sur des points à l'intérieur des cellules.
		\item Dés lors, les vecteurs de distance sont différents de $\vec{0}$ et le produit scalaire est différent de $0$. 
		\item Mais si la partie fractionnaire est proche de $0$, les points considérés sont proches des coins de cellules. Les variations sont faibles et nous avons un risque de quasi-dégénérescence. 
		\item De plus, lorsque la fréquence augmente, nous utilisons moins de points et donc moins de vecteurs de gradient. Ce qui peut entrainer des problèmes de repliement spectral (\textit{aliasing}). %+ expliquer
	\end{itemize}
\end{itemize}


